\documentclass[10pt]{article}
\usepackage[margin=0.8in]{geometry}
\usepackage{amsmath,amssymb,hyperref}
\hypersetup{colorlinks=true,urlcolor=blue,linkcolor=blue}
\setlength{\parindent}{0pt}\setlength{\parskip}{4pt}
\newcommand{\R}{\mathbb R}\newcommand{\C}{\mathbb C}
\begin{document}
\begin{center}
{\Large\bfseries Problem 00000003135: a missing connectedness hypothesis}

lizaixi01 \quad AI-led research and formalization \quad 23 October 2026
\end{center}
\textbf{Result and scope.}
The asserted bijection between the unitary dual of a nilpotent Lie group and its coadjoint orbits fails for Lie groups without a connectedness assumption. The two-element group
\[
G=\mathbb Z/2\mathbb Z
\]
with its usual discrete smooth real Lie-group structure is a counterexample. Its actual Lie algebra is zero, so there is exactly one coadjoint orbit. Its full unitary dual contains at least two different classes, represented by the trivial and sign representations. These are actual representations on the complex Hilbert space $\C$, and their inequivalence is proved for every possible unitary intertwiner.

Neither language of the source imposes connectedness or simple connectedness. We address that literal unrestricted Lie-group reading. We do not refute the connected simply connected nilpotent version of Kirillov's theorem. We also make no claim about the additional, unspecified measure assertion. The entire original English and Chinese statement is retained in \texttt{SOURCE.md}. The finite group here is an ordinary zero-dimensional \emph{real} Lie group; its real tangent Lie algebra is not a finite-field Lie ring used in a different finite-group orbit method.

\section*{1. An actual nilpotent smooth real Lie group}
The formal group is \texttt{Multiplicative (ZMod 2)}: its multiplication is addition modulo $2$. Its actual finite cardinality is $2$, and it is abelian, hence nilpotent. Give it the discrete topology. It is compact, Hausdorff and second countable.

Use the zero-dimensional real normed space $E=\R^{\varnothing}$ as the model. For each $g\in G$, the singleton $\{g\}$ is an open chart domain, mapped homeomorphically onto $E=\{0\}$. These charts give a genuine smooth manifold. Every map from a discrete manifold is locally constant and smooth. In particular multiplication $G\times G\to G$ and inversion $G\to G$ are smooth. Lean constructs the actual charted-space, manifold and Lie-group instances for every differentiability order.

We then use Mathlib's \texttt{GroupLieAlgebra}, the true tangent space at the identity with its standard Lie bracket. That tangent space is modeled on $E$, hence every tangent vector is zero. The development derives finite dimensionality and
\[
\dim_\R\mathfrak g=0.
\]
The Lie algebra structure is obtained from the existing smooth Lie-group construction, not supplied as an unrelated zero bracket.

\section*{2. The genuine coadjoint action has one orbit}
For $g\in G$, conjugation is the actual map
\[
c_g(x)=gxg^{-1}.
\]
It is smooth and, since $G$ is abelian, equals the identity. The formal adjoint map is its actual manifold derivative at the identity, with the target tangent space transported using $c_g(1)=1$:
\[
\operatorname{Ad}_g=(d c_g)_1:\mathfrak g\longrightarrow\mathfrak g.
\]
No derivative of a nonsmooth function or convention for a failed derivative is used.

On the actual dual $\mathfrak g^*=\operatorname{Hom}_\R(\mathfrak g,\R)$ define
\[
\operatorname{Ad}^*_g\ell=\ell\circ\operatorname{Ad}_{g^{-1}}.
\]
Since $\mathfrak g=0$, its dual contains exactly the zero functional. The actual quotient by the coadjoint orbit relation is therefore nonempty and a subsingleton. Thus it has one element; this is derived from the tangent space and orbit quotient, not assigned as a count.

\newpage
\section*{3. The full unitary dual, not a list of chosen characters}
The formal representation type ranges over all nonzero complex Hilbert spaces in a fixed universe. An object contains its carrier $H$, complete normed inner-product structure, a group homomorphism
\[
\rho:G\longrightarrow U(H),
\]
strong continuity, and irreducibility. Here $U(H)$ is the actual group of surjective complex-linear isometries. Irreducibility says that every closed invariant complex-linear subspace is either zero or all of $H$.

For two such representations on possibly different carriers, unitary equivalence is the existence of an actual surjective complex-linear isometry $T:H_1\to H_2$ satisfying
\[
T\rho_1(g)v=\rho_2(g)Tv\qquad(g\in G,\ v\in H_1).
\]
The proof verifies that this relation is reflexive, symmetric and transitive, and defines $\widehat G$ as the quotient of the full representation type by that relation. The carrier is not restricted to dimension one and is not replaced by Boolean labels.

Consider two objects with $H=\C$. Write $s$ for the nonidentity element of $G$:
\[
\rho_+(s)z=z,\qquad \rho_-(s)z=-z.
\]
At the identity both maps act as the identity. Their group laws are checked in the actual order-two group. Both actions consist of unitary maps and are strongly continuous because the domain is discrete. Every complex-linear subspace of $\C$ is zero or all of $\C$, proving irreducibility for both objects.

These two representations are not unitarily equivalent. If $T$ intertwined them, the equation at $s$ and $1\in\C$ would give
\[
T(1)=-T(1).
\]
Complex linearity also gives $T(-1)=-T(1)$, so injectivity would imply $-1=1$, a contradiction. The proof quantifies over all unitary intertwiners. Passing to the actual quotient therefore gives two distinct elements of $\widehat G$. Classification of all other irreducible representations is unnecessary.

\section*{4. No bijective orbit correspondence}
Let $\mathcal O=\mathfrak g^*/G$ be the coadjoint orbit quotient. It is nonempty and has at most one element, while $\widehat G$ has the two distinct classes just proved. If an equivalence $\mathcal O\simeq\widehat G$ existed, its inverse would map those classes to equal elements of $\mathcal O$; injectivity would make the classes equal. This contradiction proves
\[
\neg\exists f:\mathcal O\longrightarrow\widehat G
\quad\text{such that }f\text{ is bijective}.
\]
The necessary set-correspondence clause already fails. We do not supply a numerical value or an interpretation for the source's additional KKS-measure phrase.

\section*{5. Verification and convention reference}
The source has 51 explicit declarations: 23 theorems, 11 definitions, 7 abbreviations, 9 instances and one representation structure. The audit covers all 89 declarations in its namespace, including the structure constructor, projections, recursors and generated auxiliary declarations. The full pinned project includes a new build, a warning-as-error axiom audit, a fresh enumeration of the imported namespace, and official full fresh-kernel replay. Only the usual foundational axioms are permitted; no \texttt{sorry}, \texttt{native\_decide}, or additional axiom is used.

Fuchs and Kirillov, \emph{Representations of the group of two-diagonal triangular matrices}, Arnold Mathematical Journal (2026), Introduction, explain the connected simply connected nilpotent scope of the literal orbit-method rules and the need for corrections in other settings:
\url{https://armj.math.stonybrook.edu/Articles/250708-Fuchs/index.html}.
Only that introductory scope is used; the example and all required objects above are proved in Lean. Organizer acceptance and any controlled efficiency gain are not claimed.
\end{document}
